\documentclass[11pt]{article}
\usepackage{amsmath}
\usepackage{amssymb}
\usepackage{graphicx}
\usepackage{geometry}
\usepackage{float}
\usepackage{biblatex}
\usepackage{booktabs}
\usepackage{listings}
\usepackage{courier}
\usepackage{array}
\newcommand{\bibent}{\noindent\hangindent 20pt}
\lstset{
    basicstyle=\ttfamily,
    columns=fullflexible,
    keepspaces=true,
    breaklines=true
}

\addbibresource{reference.bib}

\geometry{
    letterpaper,
    left=24mm,
    right=24mm,
    top=24mm,
    bottom=24mm
}

\title{A Research Study on Islanders to Determine the Effect of Playing Musical Instruments on Mental Arithmetic Abilities}
\date{September 20, 2024}
\author{Boqian Mao, Jamie Tian, Junxiang Teng, Zixuan Liu}

\usepackage{titling}
\renewcommand\maketitlehooka{\null\mbox{}\vfill}
\renewcommand\maketitlehookd{\vfill\null}

\begin{document}

\begin{titlingpage}
\maketitle
\end{titlingpage}

{\newpage}
\tableofcontents

{\newpage}
\section{Abstract}
~~~~Numerous prior studies have indicated that mental exercises and sensory stimulation can enhance cognitive function, including arithmetic abilities. As a result, there is growing interest in using non-pharmacological interventions, such as music training (MT), to improve cognitive thinking. MT has been recognized for its broad cognitive benefits, affecting areas of the brain beyond the auditory and motor cortices. Research shows MT enhances neuroplasticity, improving functions such as memory, attention, and mathematical skills. These findings suggest that music-based interventions could potentially be used for educational and therapeutic purposes due to its neuroadaptive capabilities.

\section{Introduction}
~~~~There is robust neuroscience evidence suggesting that structured musical activities significantly enhance key cognitive areas crucial for academic performance, notably numerical cognition. Extensive prior research indicates that MT can substantially improve the connectivity and coordination within brain networks, which bolsters complex reasoning and problem-solving abilities. These benefits are particularly noticeable in tasks that involve counting, understanding proportions, and other fundamental mathematical concepts. Moreover, the duration and intensity of music training appear to play a vital role in determining the extent of cognitive benefits, with prolonged periods of training correlating with more substantial enhancements in cognitive function. This relationship between training duration and cognitive outcomes necessitates a detailed exploration of how various lengths of music exposure may influence specific cognitive abilities such as arithmetic skills.

Not only has music training been identified to improve cognitive function in healthy individuals, but it also holds promise as an intervention strategy for people with cognitive impairments. For example, children with developmental dyscalculia (DD) - a learning disorder characterized by significant difficulties in grasping numbers and performing mathematical calculations - may benefit from structured music training. Similarly, some studies suggest that engaging in music training during critical developmental periods such as adolescence and early adulthood could enhance cognitive reserve and delay the cognitive deterioration associated with aging (i.e. Alzheimer’s disease).

Previous research has mainly focused on the general cognitive benefits of musical training. However, there are very few studies that have examined the specific influence of different instruments and practice durations on mathematical abilities. Therefore, this research study aims to determine the effect of playing musical instruments on mental arithmetic skills. Specifically, we seek to determine how playing different musical instruments and varying the duration of practice sessions affects mathematical ability in various age groups.

In our study, it employs a randomized block design to investigate the effects of music training on arithmetic performance across different practice durations and age groups. We selected a sample of virtual islanders, dividing them into two age blocks: teenagers (13-17 years) and young adults (18-30 years). Participants were then randomized into either a control group, or one of three experimental groups. The experimental groups were assigned to practice either the piano, flute, or cello. Each participant engaged in practice sessions lasting either 10 or 20 minutes. Our response variable was the performance of each individual on arithmetic tests after the practice session. The experiment was structured into a two-way randomized block design to control for age and individual variability. Four researchers consistently collected data to ensure reliability across all treatment groups. Preliminary analysis through ANOVA suggests significant variations in arithmetic performance based on the type of instrument played and the duration of practice.

The study concludes that specific musical instruments and longer practice sessions are beneficial in enhancing arithmetic abilities, offering substantial implications for educational strategies that integrate musical training to boost mathematical skills. These findings contribute to the growing body of evidence on the cognitive benefits of MT, particularly in the field of numerical understanding. This can eventually be applied to the development of educational strategies and cognitive rehabilitation programs, aiming to harness music's potential for educational enhancement.

\section{Methods}
\subsection{Participants}

Participants were selected from the Island, a virtual human population specifically developed to facilitate learning and teaching in experimental design, and treatments were randomly assigned across the study groups. Data collection was conducted by four researchers, each responsible for an equal number of participants, to maintain consistency and control across the dataset. This structured approach ensured that each treatment group was comparable in size and composition, reinforcing the validity of the experiment.

\subsection{Design}
The study will be set up as a Two-Way Randomized Block Design. The parameters for the design are detailed here:
\begin{table}[H]
    \centering
    \begin{tabular}{|c|c|c|c|c|}
    \hline
    \multicolumn{1}{|c|}{Response Variable} & \multicolumn{4}{c|}{Arithmetic Test Correctness} \\ \hline
    \multicolumn{1}{|c|}{Treatment 1 (Instrument)} & Control & Piano & Flute & Cello \\ \hline
    Treatment 2 (Practicing Period) & \multicolumn{2}{c|}{10 minutes} & \multicolumn{2}{c|}{20 minutes} \\ \hline
    Blocking (Age Group) & \multicolumn{2}{c|}{13-17} & \multicolumn{2}{c|}{18-30} \\ \hline
    \end{tabular}
    \caption{Experimental Design}
\end{table}


The factor diagram is detailed below:
\begin{figure}[htt!]
    \centering
    \includegraphics[width=0.7\linewidth]{Factor Diagram.png}
    \caption{Factor Diagram}
\end{figure}

We selected three distinct types of musical instruments—Piano, Cello, and Flute—to explore their differential effects on the enhancement of arithmetic abilities. The study also examines the impact of practice duration by incorporating two levels of playing time: 10 minutes and 20 minutes, to assess whether a longer practice period influences arithmetic performance. To account for inherent differences in cognitive abilities, participants were stratified into two age blocks: teenagers (13-17 years) and young adults (18-30 years). This blocking approach ensures that the age-related variability in arithmetic skills is controlled throughout the experiment.

\subsection{Instruments}

Arithmetic abilities were initially assessed using the Mental Arithmetic Basic Test, which consists of 40 arithmetic problems; initial scores were recorded as the number of correct responses out of 40. Participants in the treatment groups then practiced their assigned instruments—Piano, Cello, or Flute—for either 10 or 20 minutes. The control group underwent a waiting period of equivalent duration between the two tests without engaging in any activities, allowing for a controlled comparison of the effects of musical practice versus inactivity on arithmetic performance changes.


\subsection{Procedure}

\textbf{Step 1:} Determine the number of participants for each treatment. Recruit subjects from the Island, ensuring an equal distribution between ages 13-17 and 18-30. Randomly assign subjects to specific treatments, maintaining balance across age groups.\\

\noindent \textbf{Step 2:} Randomly allocate the groups, already divided by age blocks, to different treatment conditions as follows: \begin{itemize} \item Control group: No activity for 10 minutes, no activity for 20 minutes. \item Piano group: Playing piano for 10 minutes, playing piano for 20 minutes. \item Flute group: Playing flute for 10 minutes, playing flute for 20 minutes. \item Cello group: Playing cello for 10 minutes, playing cello for 20 minutes. \end{itemize}\\

\noindent \textbf{Step 3:} Administer the initial Mental Arithmetic Basic Test to each participant to record the baseline arithmetic performance.\\

\noindent \textbf{Step 4:} Implement the designated treatments—participants engage in playing the assigned musical instrument or observe the specified waiting period.\\

\noindent \textbf{Step 5:} Post-treatment, participants retake the Mental Arithmetic Basic Test to assess any changes in arithmetic performance.\\

\noindent \textbf{Step 6:} Calculate the difference in scores on the Mental Arithmetic Basic Test before and after treatment for each participant. This difference serves as the response variable for subsequent analysis.\\



\newpage
\section{Data Analysis}
\subsection{Type of Statistical Analysis}
We will conduct a ANOVA to analyze the data collected from our experiment. The primary goal is to examine whether different musical instruments (Piano, Cello, Flute) and practice durations (10 minutes, 20 minutes) have a significant effect on arithmetic test performance, and whether there is an interaction between instrument type and practice duration. The analysis will be performed using the R statistical software. We will use the aov() function to model the effects of these factors and their interaction, while also including age as a blocking factor to control for variability due to participant age (13-17, 18-30). Post-hoc comparisons will be conducted using Tukey’s HSD test to explore specific group differences where significant main or interaction effects are found. Residuals will be checked for normality and homogeneity of variance using diagnostic plots to validate the assumptions of ANOVA. 

\subsection{Sample size determination}
We determined the necessary sample size for our experiment using G*Power with a power of 0.80, giving us an 80\% probability of correctly rejecting the null hypothesis if there is a true effect. We set the Alpha level at 0.05, meaning there's a 5\% chance of a Type I error. For the effect size, we chose a moderate value of 0.25, which represents a reasonable expected difference between the groups. Our experiment consists of 16 groups based on the combinations of musical instruments and practice durations, with one covariate (age group) to account for age-related variability. After running the analysis, G-Power calculated that a sample size greater than 179 will be appropriate for this study. We finally determined a sample size of 192 for our study, ensuring adequate representation in each group. 
\begin{figure}[h]
    \centering
    \includegraphics[width=0.7\textwidth]{diagram}
    \caption{G-Power Sample Size Calculation}
    \label{fig:diagram}
\end{figure}


\newpage
\section{Results}

\subsection{Box Plots}
In each box plot, the black bar in the middle of each box represents the median of difference in mental arithmetic scores. For each box, the maximum height represents the third quartile and the minimum height represents the first quartile. The black dots represents potential outliers. 
\begin{figure}[H]
    \centering
    \includegraphics[width=1\linewidth]{boxplots.pdf}
    \caption{Box Plot Comparing in Mental Arithmetic Scores between Age Group, Type of Musical Instrument, and Practicing Period. }
\end{figure} 
The differences among the medians, first quartiles, and third quartiles of all the box plots suggest a significant change in mental arithmetic scores of different types of musical instrument. There is no significant difference between the age groups and practising period. There are some potential outliers that need further investigation. 

\begin{figure}[H]
    \centering
    \includegraphics[width=1\linewidth]{boxplot.pdf}
    \caption{Box Plot of Treatment 1 and Treatment 2 with Change in Mental Arithmetic Scores. }
    \label{fig:enter-label}
\end{figure}
The difference among the medians, first quartiles, and third quartiles between the two box plots for each type of musical instrument suggests that there is a significant change in mental arithmetic scores between the practicing periods for Cello. There is no significant difference between Flute and Piano. A few potential outliers need to be investigated. 


\subsection{Interaction Plots}
\begin{figure}[H]
    \centering
    \includegraphics[width=1\linewidth]{interactions.pdf}
    \caption{Interaction Plot of Treatment 1 and Treatment 2 with Change in Mental Arithmetic Scores. }
    \label{fig:enter-label}
\end{figure}
The interaction plot shows non-parallel lines, which suggest that there is a significant interaction between the type of musical instrument and practicing period. There is an interaction between the cello and flute treatments.

\subsection{Residual Diagnositics}
\begin{figure}[H]
    \centering
    \includegraphics[width=1\linewidth]{residuals.pdf}
    \caption{Residuals Plots for ANOVA model. }
    \label{fig:enter-label}
\end{figure}
In the Residual v.s. Fitted plot, there is a light curve pattern, which suggests that there may be some non-linearity in the model, and a pattern that is not constant in the spread of residuals, which suggests heterscedasticity. In the Normal Q-Q plot, the points can be fitted with a straight line mostly but there are heavy tails and points far from the line, which respectively suggest that normality is violated and that there are potential outliers. In the Scale-Location plot, there is some changes in the spread of the variance, which suggests heterscedasticity. The Residuals v.s. Factor Levels plot suggests a few influential points that need further investigation.

\subsection{Tukey HSD Adjusted P-values}
\begin{table}[H]
  \centering
    \begin{tabular}{lccccc}
        \toprule
        Comparison &Difference &Lower &Upper &P value Adjusted
        \\
        \midrule
        Control-Cello& -3.875000& -5.5826206& -2.1673794& 0.0000001
        \\
        Flute-Cello& -1.208333& -2.9159539& 0.4992872& 0.2604706
        \\
        Piano-Cello& -2.375000& -4.0826206& -0.6673794& 0.0022493
        \\
        Flute-Control& 2.666667& 0.9590461& 4.3742872& 0.0004380
        \\
        Piano-Control& 1.500000& -0.2076206& 3.2076206& 0.1070813
        \\
        Piano-Flute& -1.166667& -2.8742872& 0.5409539& 0.2904521
        \\
        \bottomrule
    \end{tabular}
  \label{tab:grav1}
  \caption{Post-Hoc Analysis of Differences in Change of Mental Arithmetic Scores between Different Musical Instrument Treatment Types.}
\end{table}
P-values of 0.0000001, 0.0022493, and 0.0004380, which are less than a critical value of 0.05, respectively suggest that the changes in mental arithmetic scores between control and cello treatments, piano and cello treatments, and flute and control treatments are statistically significantly different. P-values of 0.2604706, 0.1070813, and 0.2904521, which are greater than a critical value of 0.05, respectively suggest that the changes in mental arithmetic scores between flute and cello treatments, piano and control treatments, and piano and flute treatments are not statistically significantly different.

\subsection{ANOVA Analysis}
\begin{table}[H]
  \centering
    \begin{tabular}{lccccc}
        \toprule
        &Df &Sum Sq &Mean Sq &F value   &Pr(\(>\)F)  
        \\
        \midrule
        Treatment\textunderscore 1 &3 &394.1 &131.35 &12.618 &1.55e-07
        \\
        Treatment\textunderscore 2 &1 &8.3 &8.33 &0.800 &0.3721
        \\
        Age\textunderscore Group &1 &7.5 &7.52 &0.722 &0.3965
        \\
        Treatment\textunderscore 1:Treatment\textunderscore 2 &3 &119.5 &39.83 &3.826 &0.0109
        \\
        Residuals &183 &1905.1 &10.41
        \\
        \bottomrule
    \end{tabular}
  \label{tab:grav1}
  \caption{Two-way ANOVA table with Blocking and Interactions.}
\end{table}
P-values of \( 1.55 \times 10^{-7} \) and 0.0109, which are less than a critical value of 0.05, respectively suggest that the type of musical instrument and the interaction between the type of musical instrument and practicing period have statistically significant effects on the change of mental arithmetic test scores. P-values of 0.3721 and 0.3985, which are greater than a critical value of 0.05, respectively suggest that the practicing period and age group do not statistically significantly affect the change of mental arithmetic test scores. A p-value of 0.37785 also suggests an ineffective age group blocking.

\section{Discussion}
\subsection*{Objective Recap}
The objective of this study is to investigate the impact of different musical instrument types, age groups, and practicing periods on mental arithmetic test scores. By analyzing interactions between these factors, the study aims to determine how the combination of musical training variables affects cognitive performance.

\subsection*{Key Findings}
Our experiment employs a two-way ANOVA analysis with blocking on age groups and musical instruments, using a sample size of 192 participants. We test for the interaction between treatment groups, including musical instrument types (cello, flute, piano, and control), practicing periods, and age groups, with a critical alpha level of 0.05 to determine statistical significance. Post-hoc Tukey HSD analysis is also conducted to explore pairwise differences between treatments.

\noindent{From the ANOVA test, we have evidence that the type of musical instrument and the interaction between the type of musical instrument and practicing period statistically affect mental arithmetic test score. However, we do not have sufficient evidence that the practicing period and age group statistically affect the mental arithmetic test score, indicating that the individual factors may be less influential without considering interactions. We do not have evidence that blocking by age group is effective as well.}

\noindent{From Tukey HSD post-hoc comparison results, we have evidence that suggests each of the following two treatments, control and cello treatments, piano and cello treatments, and flute and control treatments, affect the difference in mental arithmetic scores between each two of them. However, flute and cello treatments, piano and control treatments, and piano and flute treatments, do not significantly affect the differences in mental arithmetic scores between each pair of them.}

\subsection*{Box Plot and Residual Analysis}
The box plots show that age, musical instrument type, and practice period all have an effect on mental arithmetic scores. A few potential outliers require further investigation. The residual diagnostic plots show non-constant variance in the distribution of residuals, showing heteroscedasticity. Furthermore, thick tails in the Q-Q plot indicate that normality assumptions were marginally broken, while the Scale-Location plot demonstrates variability across the sample. These diagnostic difficulties indicate that, while the model is essentially correct, potential outliers and heteroscedasticity warrant further research. Interaction plots show non-parallel lines across treatments, indicating strong interactions, especially between the cello and flute treatments. 


\subsection*{Limitations:}
Several limitations in the study design could have influenced our findings. First, the observed heteroscedasticity and probable outliers indicate that the model may not fully reflect the complexities of the relationships between musical instrument type and cognitive progress. Furthermore, the absence of control over external factors, such as participants' previous experience with mental arithmetic, may have caused variability. Future research should seek to adjust for these external variables in order to better isolate the effects of musical instruction.

\noindent{To improve the reliability of these findings, future studies should use a bigger sample size and account for individuals' background in both music and cognitive activities. Additional elements, such as the duration of musical instruction or specific cognitive exercises other than mental arithmetic, should be investigated to acquire a better understanding of how musical practice impacts various forms of cognitive function.}

\subsection*{Conclusion:}
This study provides compelling evidence that musical training, particularly with instruments like the cello and piano, has a measurable impact on mental arithmetic performance. The interactions between musical instrument types and practicing periods significantly affect cognitive outcomes, as demonstrated by the statistically significant results in our ANOVA analysis. These findings suggest that musical training could be an effective method to enhance cognitive function, though further research is needed to refine these conclusions.

\newpage
\section{Reference}
\bibent
Ciurea, Alexandru Vlad, et al. “Cognitive Crescendo: How Music Shapes the Brain’s Structure and Function.” Brain Sciences, 13(10), 2023, www.mdpi.com/doi/10.3390/brainsci13101390.
\vspace{10pt}

\sloppy
\bibent 
Gaab, Nadine, and Jennifer Zuk. “Is There a Link between Music and Math?” \textit{Scientific American}, Scientific American, 20 Feb. 2024, www.scientificamerican.com/article/is-there-a-link-between-music-and-math/.
\vspace{10pt}

\bibent
Hudak, E. M., et al. “Keys to Staying Sharp: A Randomized Clinical Trial of Piano Training Among Older Adults with and without Mild Cognitive Impairment.” Contemporary Clinical Trials, 84, 2019, doi.org/10.1016/j.cct.2019.06.003.
\vspace{10pt}

\bibent
Husebø, Bettina Sandgathe, et al. “The Effect of Music-Based Intervention on General Cognitive and Executive Functions, and Episodic Memory in People with Mild Cognitive Impairment and Dementia: A Systematic Review and Meta-Analysis.” Healthcare, 10(8), 2022, doi.org/10.3390/healthcare10081462.
\vspace{10pt}

\bibent 
Jiménez-Palomares, María, et al. “Benefits of Music Therapy in the Cognitive Impairments of Alzheimer’s-Type Dementia: A Systematic Review.” Journal of Clinical Medicine, U.S. National Library of Medicine, 1 Apr. 2024, www.ncbi.nlm.nih.gov/pmc/articles/PMC11012733/.
\vspace{10pt}

\bibent 
Miendlarzewska, Ewa A, and Wiebke J Trost. “How Musical Training Affects Cognitive Development: Rhythm, Reward and Other Modulating Variables.” \textit{Frontiers in Neuroscience}, U.S. National Library of Medicine, 20 Jan. 2014, www.ncbi.nlm.nih.gov/pmc/articles/PMC3957486/.
\vspace{10pt}

\bibent 
Ribeiro, Fabiana Silva, and Flávia Heloísa Santos. “Persistent Effects of Musical Training on Mathematical Skills of Children with Developmental Dyscalculia.” \textit{Frontiers in Psychology}, U.S. National Library of Medicine, 10 Jan. 2020, www.ncbi.nlm.nih.gov/pmc/articles/PMC6965363/.
\vspace{10pt}

\bibent 
Román-Caballero, R., et al. “Please Don’t Stop the Music: A Meta-Analysis of the Cognitive and Academic Benefits of Instrumental Musical Training in Childhood and Adolescence.” *Educational Research Review*, 35, 2022, doi.org/10.1016/j.edurev.2022.100436.
\vspace{10pt}

\bibent
Sihvonen, A. J., et al. “Effects of a 10-week Musical Instrument Training on Cognitive Function in Healthy Older Adults: Implications for Desirable Tests and Period of Training.” *Frontiers in Aging Neuroscience*, 15:1180259, 2023, doi.org/10.3389/fnagi.2023.1180259.
\vspace{10pt}

\bibent
Sousa, J., et al. “When Music and Long-Term Memory Interact: Effects of Musical Expertise on Functional and Structural Plasticity in the Hippocampus.” *PLOS ONE*, 15, 2020, doi.org/10.1371/journal.pone.0013225.
\vspace{10pt}

\bibent
Verghese, J., et al. “Leisure Activities and the Risk of Dementia in the Elderly.” *New England Journal of Medicine*, 348, 2003, doi.org/10.1056/NEJMoa022252.
\vspace{10pt}

\bibent 
Zhang, Lei, et al. “Musical Experience Offsets Age-Related Decline in Understanding Speech-in-Noise: Type of Training Does Not Matter, Working Memory Is the Key.” \textit{Ear and Hearing}, U.S. National Library of Medicine, 2021, www.ncbi.nlm.nih.gov/pmc/articles/PMC7969154/.
\vspace{10pt}
 
% \printbibliography

\end{document}
